% M1: one document per student. A filled copy of template-m1.tex.
% Build: latexmk -pdf m1.tex   (pdflatex + biber)
\documentclass[11pt,a4paper]{article}
\usepackage[margin=2.5cm]{geometry}
\usepackage[T1]{fontenc}
\usepackage[utf8]{inputenc}
% TeX Gyre Heros, a Helvetica clone with the same letter widths as Arial in the Word template.
\usepackage{tgheros}
\renewcommand\familydefault{\sfdefault}
\usepackage{xcolor}
\usepackage{enumitem}
\usepackage{array}
\usepackage{graphicx}
% The Ensemble diagram in section 2.2.
\usepackage{tikz}
\usetikzlibrary{fit,positioning}
% Styles for Figures 1 and 2 (section 2.1). The Ensemble figure sets its own box and lab styles.
\tikzset{
  rdfig/.style={font=\footnotesize, >=stealth},
  box/.style={draw, minimum height=5mm, inner sep=2.5pt, align=center},
  lm/.style={box, fill=black!12},
  lab/.style={font=\scriptsize, inner sep=1.5pt, align=center},
}
\usepackage[british]{babel}
\usepackage[style=american]{csquotes}
\usepackage[style=ieee,backend=bibtex]{biblatex}
% The reference list follows IEEE punctuation and dates ("A, B, and C", "Nov. 9, 2026").
\DeclareLanguageMapping{british}{american}
\addbibresource{references.bib}
\addbibresource{course.bib}
\addbibresource{merkl.bib}
\usepackage[hidelinks]{hyperref}
\hypersetup{pdftitle={M1: Literature survey, agent design and project plan}}
\setlength\parindent{0pt}
\setlength\parskip{0.5em}
% Less white space around the one figure and its caption.
\setlength\intextsep{6pt}
\setlength\abovecaptionskip{4pt}
\setlength\belowcaptionskip{0pt}

% The hints, and the example rows of the table, in grey italic.
\newif\ifhints
\hintsfalse   % hints are hidden
\usepackage{comment}
\ifhints
  \newenvironment{hint}{\par\color{gray}\itshape\setlist{nosep,leftmargin=1.5em}}{\par}
  \newcommand\examplerow[3]{{\color{gray}\itshape #1} & {\color{gray}\itshape #2} & {\color{gray}\itshape #3} \\}
\else
  \excludecomment{hint}
  \newcommand\examplerow[3]{}
\fi

% ---------------------------------------------------------------------------
\newcommand\TODO[1]{\textcolor{red}{#1}}
\newcommand\myname{Leonie Merkl}
\newcommand\myteam{A Team}
\newcommand\Imp{Jurcik}     % Improviser
\newcommand\Mus{Merkl}   % Musicologist
\newcommand\Pla{Mukherjee}         % Player
% ---------------------------------------------------------------------------

\newcommand\code[1]{\texttt{#1}}
% A clickable reference to a section, printed as §1, §2.1, ...
\newcommand\secref[1]{\hyperref[#1]{\S\ref*{#1}}}
% A paragraph that starts with a bold lead-in. The bold text marks the new paragraph,
% so the space between paragraphs (\parskip) is taken back above it.
\newcommand\lead[1]{\vspace{-\parskip}\textbf{#1}}

\begin{document}

{\Large\bfseries M1 --- Literature survey, agent design and project plan}\par\medskip
\textbf{Name:} \myname \quad \textbf{Team:} \myteam \quad \textbf{Specialisation:} Musicologist
(Research)\par\medskip

\section{Literature survey}\label{sec:survey}

% >>> OWN PART: literature survey <<<
\lead{Retrieval.} Retrieval-augmented generation (RAG) first retrieves passages for the input and then
generates the answer from them. The model itself does not decide whether or when to retrieve
\parencite{lewis2020retrieval}. Its dense retriever, DPR, was better than BM25 on all datasets tested except SQuAD
\parencite{karpukhin2020dense}. At the Reading Desk, only the titles can be searched for free, so retrieval starts from them.

\lead{Agentic reading.} ReAct alternates reasoning and tool calls. Non-informative searches caused 23\% of its
errors, and it often repeated earlier thoughts and actions \parencite{yao2023react}. Unlike RAG that retrieves
once or follows a predefined workflow, A-RAG lets the model choose among keyword search, semantic search and
chunk reading, and it was more accurate than traditional RAG with comparable or fewer retrieved tokens
\parencite{du2026arag}. DeepRead shows the model the document's headings and lets it choose between searching
and reading. It mostly searched first and then read \parencite{li2026deepread}. At the Reading Desk, each read is charged in characters as well as tokens.

\lead{Cost and context.} In Wu et al., each piece of evidence has a price and each question a budget to spend on it. No fixed selector
uniformly dominated, and one model often consumed its budget instead of managing it
\parencite{wu2026knowledge}. Joren et al.\ found that raising the context limit from 6,000 to 10,000 tokens
left the share of questions with sufficient context effectively unchanged \parencite{joren2025sufficient}. The Reading Desk sets no budget itself and only counts the cost against the score.

\lead{Evaluation.} Kapoor et al.\ argue that agents should be judged on accuracy and cost together, as a Pareto
front: on HumanEval, simple baselines matched or beat complex agents at much lower cost
\parencite{kapoor2024ai}. ALCE tests whether citations jointly support each statement and whether any citation
is irrelevant \parencite{gao2023enabling}.

In a workflow, models and tools follow predefined code paths, whereas in an agent the model directs its own
process and tool use \parencite{anthropic2024building}. By this definition, classic RAG is a workflow, while
A-RAG and DeepRead are agents. The Reading Desk combines agentic reading, a charge per read and a score on the exact supporting set. No source in this survey treats all three together, so the design draws on each strand: term matching, agentic reading, a budget in code and cost-aware evaluation.

\section{Agent design}

\subsection{My agent}\label{sec:agent}

% >>> OWN PART: my agent (four questions), Figures 1 and 2 <<<
\lead{The problem.} Only \code{browse}, which lists titles and lengths, is free. \code{examine} (abstract and
section headings) and \code{read} (one section) cost characters, even when repeated. A right answer scores 1, a wrong one or a skip 0. The named volumes score 1 if they are exactly a set that establishes
the answer. An extra volume takes off a share, and a missing one gives 0. Extra reading and tokens also count,
with weights published before M2 \parencite{course-api}. If the weights allow, no question
will be skipped, and one volume too many will be named rather than one too few.

\lead{Approach.} The \emph{plain-prompt baseline} gives the model only the rules and the arena's messages. It
plays alone with the arena's queries, without help from code and with a fixed limit per question
\parencite{course-s4}. It is the reference point and will be built first. The agent then adds classical
computation in code: \code{bm25\_search}, \code{find\_in\_toc} and regex. BM25 costs nothing, and term matching
is sensitive to highly selective keywords \parencite{karpukhin2020dense}, which names, titles of works and
years are expected to be. Dense retrieval is left out for now because titles are short, but it could be added
later.

The \emph{agentic approach} (Fig.~\ref{fig:agentic}) is the main design. The titles from \code{browse} are
in the prompt, and the model chooses every tool in a ReAct loop \parencite{yao2023react}. As in A-RAG, no
fixed search step comes before its first decision \parencite{du2026arag}. The model can therefore react to
what it finds: it can open a second volume when one is not enough, or search again when
BM25 misses a synonym. Because each step costs tokens, the code guards the loop against the repetition reported for
ReAct \parencite{yao2023react}. As seen in Wu et al., a model may use up its budget instead of managing it
\parencite{wu2026knowledge}, so the budget is enforced in code. When the loop stops, the best guess
is sent instead of a skip \parencite{course-api}.

\begin{figure}[!ht]\centering
\begin{tikzpicture}[rdfig]
  \node[box, minimum width=15mm] (q) at (0.2,0) {Question};
  \node[box, minimum width=27mm] (bm) at (3.1,0) {Prompt: question,\\titles from \code{browse}};
  \node[lm,  minimum width=34mm] (lm) at (7.4,0) {Model chooses\\the next tool};
  \node[box, minimum width=27mm] (ans) at (12.6,0) {\code{answer} with\\supporting volumes};
  \node[box, minimum width=44mm] (tool) at (4.0,-1.7) {Tools: \code{bm25\_search}, \code{find\_in\_toc},\\\code{examine}, \code{read}, regex};
  \node[box, minimum width=44mm] (grd) at (10.7,-1.7) {Guards: budget, repeated\\action, no progress};
  \draw[->] (q) -- (bm);
  \draw[->] (bm) -- (lm);
  \draw[->] (lm) -- node[lab, above] {answer} (ans);
  \draw[->] (5.95,0 |- lm.south) -- node[lab, left] {action} (5.95,0 |- tool.north);
  \draw[->] (tool) -- node[lab, above] {observation} (grd);
  \draw[->] (8.8,0 |- grd.north) -- node[lab, right] {continue} (8.8,0 |- lm.south);
  \draw[->] (grd.north -| ans.south) -- node[lab, right] {stop:\\best guess} (ans.south);
\end{tikzpicture}
\caption{Agentic approach: the model decides each step; code runs tools and guards the loop.}\label{fig:agentic}
\end{figure}

The \emph{workflow approach} is the alternative considered. Like classic RAG (\S1), it follows a path set by
code: it searches the titles, examines the top volumes, reads sections and calls the model only when
checks in code are unclear. It would spend fewer tokens, but was not chosen because a fixed path cannot
react when a question needs a second volume or a title uses a synonym.

\lead{How it will be judged.} Correct answers, characters read and tokens will be measured, and the named
volumes will be checked by citation recall and precision \parencite{gao2023enabling}. Runs will use questions
kept out of tuning, several seeds and a fixed temperature \parencite{hochlehnert2025sober,course-s3}. The agent
works if it answers more questions correctly than the baseline within the budget per question. This budget
starts at about 6,000 tokens \parencite{joren2025sufficient}, estimated here as 24,000 characters. The Pareto
front of cost and correct answers will then be used to choose between settings \parencite{kapoor2024ai}. If the
agent falls short, each wrong answer will be traced to its cause: either the right section was never read, or
it was read and the answer was still wrong. The more frequent cause will be fixed first.

\lead{Where it meets the team.} The Musicologist gets a question and a budget from the Conductor and returns an answer,
supporting volumes, a confidence and the cost, which the Conductor checks. Its model calls go through the
shared wrapper. The shared store keeps the Musicologist's questions, the sections it has read and its best answer so
far. For a new question, the sections already read are searched first, as a repeated read is charged again
\parencite{course-api}. If the Musicologist fails, the Conductor can send the stored answer instead of a skip. In the Overture, the Musicologist
rates signs and volumes by BM25 against open questions, without tokens (\secref{sec:ensemble}).

\subsection{The Ensemble}\label{sec:ensemble}

\textbf{Dependencies.} The parts depend on each other in the four ways named by
\textcite{malone1994interdisciplinary}. There are two \emph{shared resources}: the places on the trolley and
the one connection to the arena. As \emph{producer and consumer}, the Improviser delivers the volumes that the
Musicologist can read only in the hall, and a door opened by the Player gives the Improviser a shortcut.
\emph{Simultaneity} is needed at a door, where a puzzle can be answered only while the trolley stands there
\parencite{course-api}. Each question is a \emph{task with subtasks}: choosing rooms, fetching, reading and
answering.

\begin{figure}[!ht]\centering
\begin{tikzpicture}[font=\footnotesize, >=stealth,
  box/.style={draw, minimum height=5mm, inner sep=2.5pt, align=center},
  lab/.style={font=\scriptsize, inner sep=1.5pt, align=center}]
  % Outside the container, on the left: the arena (the game) and the model endpoint (the language model).
  % They are two servers: model calls never pass through the arena.
  \node[box, minimum width=14mm, anchor=east] (arena) at (0.35,0.55) {Arena};
  \node[box, minimum width=14mm, anchor=east] (llm) at (0.35,-0.55) {Model endpoint};
  % Inside: the Conductor above, the wrapper and the three specialists below, the store on the right.
  \node[box, minimum width=87mm] (con) at (6.1,0.55) {\textbf{Conductor}: loop, routing rule, checks, log};
  \node[box, minimum width=22mm] (wrap) at (2.85,-0.55) {Model wrapper};
  \node[box, minimum width=17mm] (imp) at (5.85,-0.55) {Improviser};
  \node[box, minimum width=20mm] (mus) at (7.85,-0.55) {Musicologist};
  \node[box, minimum width=14mm] (pla) at (9.7,-0.55) {Player};
  \node[draw, densely dotted, fit=(imp)(mus)(pla), inner sep=2pt] (specs) {};
  \node[box, minimum width=20mm, minimum height=16mm] (store) at (12.45,0) {Shared\\store};
  \node[lab] at (4.2,0) {decision $\downarrow$ proposal $\uparrow$};
  \node[draw, dashed, fit=(con)(wrap)(specs)(store), inner sep=3pt] (box) {};
  \draw[<->] (arena) -- node[lab, below=1pt, pos=0.4] {one\\connection} (con);
  \foreach \n in {imp,mus,pla} \draw[<->] (\n.north) -- (\n.north |- con.south);
  % Every specialist may call the wrapper, and only the wrapper calls the model.
  \draw[->] (specs.west) -- node[lab, below] {call} (wrap.east);
  \draw[->] (wrap) -- (llm);
  % The Conductor and every specialist read and write the store.
  \draw[<->] (con.east) -- (con.east -| store.west);
  \draw[<->] (specs.east) -- (specs.east -| store.west);
\end{tikzpicture}
\caption{The container (dashed) with its three specialists (dotted). The Conductor and the specialists use
the shared store. Only the specialists call the model, through the wrapper.}\label{fig:ensemble}
\end{figure}

\lead{Coordination.} The Ensemble is organised as a hierarchy \parencite{course-s2}. At the top is the Conductor, written in plain code. It holds the connection, runs the loop, checks every action and
hands each decision to one specialist (Fig.~\ref{fig:ensemble}). It also keeps the shared store: the setlist texts, sent only once, the map, the texts read and the costs. In a base quest, the first message names the quest and so the specialist. In the Overture the order of work is open, so a fixed rule settles where the trolley
goes next: the Musicologist rates each candidate sign or volume, the Improviser gives its distance, and the
Conductor takes the candidate with the most worth per distance.

The alternative considered was a blackboard \parencite{erman1980hearsay}, on which a control component selects which part writes next \parencite{course-s2}. Its authors noted that such control \enquote{takes time and requires many fairly sophisticated mechanisms} and is justified where no explicit algorithm exists \parencite[246]{erman1980hearsay}. Here a simple rule chooses the next step without using tokens, so only the shared store was taken over.

\lead{When a part fails, is slow or is wrong.} If a specialist fails, the Conductor falls back on a default without tokens. For a question, it sends the best answer so far, which the Musicologist keeps up to date in the shared store. A puzzle is skipped, and on the map the nearest exit is taken.
 A slow model call is stopped before the endpoint's 300-second cut-off \parencite{course-s4} and retried once, so a turn stays below the ten minutes of silence that end a run \parencite{course-api}. Against wrong answers, every proposal is first checked in code, as language models were found no better at verifying solutions than at producing them \parencite{kambhampati2024position}. A failed proposal is returned once with the reason. If it fails again, the last well-formed one is sent, as a wrong answer costs no more than a skip \parencite{course-api}, except on claims, which a team of three does not get.

\lead{Logging.} The starter code logs every arena message and model call \parencite{course-s4}. For each action, the log will add the deciding part, its reason, its tokens, the change in \code{spent} and the check result.

\section{Project plan}

\subsection{The team's plan}

\begin{tabular}{@{}>{\raggedright\arraybackslash}p{0.14\linewidth}>{\raggedright\arraybackslash}p{0.62\linewidth}>{\raggedright\arraybackslash}p{0.18\linewidth}@{}}
\hline
\textbf{By} & \textbf{What will be working} & \textbf{Who} \\
\hline
16 Nov & Self-test passes, and fails when the prompt is broken & All \\
20 Nov & Conductor shares: loop and routing (\Imp), shared store (\Mus), logging and failure rules (\Pla) & one each \\
23 Nov & Each base quest plays a practice run via the Conductor & each owner \\
27 Nov & \textbf{First Overture run with all parts in one container} & All \\
30 Nov week & Rehearsal: image connects to course arena and endpoint & \Imp \\
4 Dec & Overture end to end with failure rules, baselines measured & \Pla \\
8 Dec & Image frozen, all dependencies inside, reflections drafted & \Mus \\
10 Dec & M2 handed in & each \\
\hline
\end{tabular}

Each member owns one Conductor share. Open rows are checked weekly and work is moved.

\subsection{My plan}

% >>> OWN PART: my plan <<<
\begin{tabular}{@{}>{\raggedright\arraybackslash}p{0.14\linewidth}>{\raggedright\arraybackslash}p{0.83\linewidth}@{}}
\hline
\textbf{By} & \textbf{What will be working} \\
\hline
20 Nov & Shared store with the best-answer field, baseline runs \\
23 Nov & Tools, practice run via the Conductor \\
27 Nov & Agentic approach with guards, ratings of signs and volumes \\
1 Dec & Measuring script, comparison with the baseline on held-out questions \\
7 Dec & Tuning along the Pareto front, settings fixed \\
\hline
\end{tabular}

Tuning is the step most likely to slip, and it would then shrink to a small grid search.

\section*{References}

{\sloppy\printbibliography[heading=none]}

\subsection*{Use of AI}
% >>> OWN PART: AI note <<<
Claude Opus 5.5 and Claude Fable 5.2 (Anthropic) were used to summarise papers, to work out the design in conversation, to draft and edit the
text and to produce the process diagram. The papers were read by the author, and every source, quotation and
number cited from them was checked against the original. The design choices are the author's own, and the
final text and figure were checked against the author's notes and the API reference.

\end{document}
